% ============================================================
% content/solutions.tex —— 2018 全国Ⅰ卷（理综物理部分）· 详细解析
% 解析类型：source_embedded（转录原卷官方解析，公式逐条复核）
% ============================================================

\gkpart{参考答案与详细解析}

\begin{answer}
\gksolhead{14}{单项选择题}
\gkanswer{答案 B}
列车做初速度为零的匀加速直线运动，加速度为 $a$，则经过时间 $t$ 的速度 $v=at$、位移 $s=\frac{1}{2}at^{2}$，动能
\[
  E_{\mathrm{k}}=\frac{1}{2}mv^{2}=\frac{1}{2}ma^{2}t^{2},
\]
即动能与时间的平方成正比，A 错误。又
\[
  E_{\mathrm{k}}=\frac{1}{2}mv^{2}=m\cdot\frac{v^{2}}{2}=ma\cdot\frac{1}{2}at^{2}=mas,
\]
故动能与位移成正比，B 正确。动能与速度的平方成正比（C 错误）、与动量的平方成正比（D 错误）。

故选 B。
\end{answer}

\begin{answer}
\gksolhead{15}{单项选择题}
\gkanswer{答案 A}
设物块质量为 $m$、弹簧劲度系数为 $k$，静止时弹簧压缩量为 $x_{0}$，有 $kx_{0}=mg$。

施加竖直向上的力 $F$ 后，物块做匀加速直线运动，加速度为 $a$（恒定）。当物块离开静止位置的位移为 $x$ 时，弹簧压缩量为 $x_{0}-x$，由牛顿第二定律
\[
  F+k(x_{0}-x)-mg=ma .
\]
代入 $kx_{0}=mg$ 得
\[
  F=ma+kx .
\]
可见 $F$ 与 $x$ 成一次函数关系，且 $x=0$ 时 $F=ma>0$（截距为正），斜率 $k>0$。在弹簧恢复原长前（$0\le x\le x_{0}$）该关系始终成立。

对应图像应为一条纵截距为正、斜率为正的直线，即 A 图。

故选 A。
\end{answer}

\begin{answer}
\gksolhead{16}{单项选择题}
\gkanswer{答案 D}
由 $ab=5\,\si{\centi\metre}$、$bc=3\,\si{\centi\metre}$、$ca=4\,\si{\centi\metre}$ 知 $ab^{2}=bc^{2}+ca^{2}$，故 $\triangle abc$ 为直角三角形，$\angle c=90^{\circ}$，且 $c$ 到 $ab$ 的垂足即 $c$ 在 $ab$ 上的投影点。

小球 $c$ 所受合力沿 $a$、$b$ 的连线方向，说明 $a$、$b$ 对 $c$ 的库仑力在垂直于 $ab$ 方向的分量互相抵消。设 $c$ 到 $ab$ 的垂距为 $h_{c}$，$a$、$b$ 到垂足的距离分别为 $d_{a}$、$d_{b}$，则
\[
  d_{a}=\frac{ac^{2}}{ab}=\frac{16}{5}\,\si{\centi\metre},\qquad
  d_{b}=\frac{bc^{2}}{ab}=\frac{9}{5}\,\si{\centi\metre},\qquad
  h_{c}=\frac{ac\cdot bc}{ab}=\frac{12}{5}\,\si{\centi\metre}.
\]
两力垂直于 $ab$ 的分量同向相加、不能抵消，只有在两力沿垂线方向分量**反向**（即 $a$、$b$ 带异种电荷）时才可能抵消，故 $a$、$b$ 必异号，A、C 错误。

由垂直分量抵消：
\[
  k\frac{|Q_{a}|Q_{c}}{ac^{2}}\cdot\frac{h_{c}}{ac}
  =k\frac{|Q_{b}|Q_{c}}{bc^{2}}\cdot\frac{h_{c}}{bc},
\]
化简得
\[
  \frac{|Q_{a}|}{|Q_{b}|}=\frac{ac^{3}}{bc^{3}}=\frac{4^{3}}{3^{3}}=\frac{64}{27},
\]
即 $k=\dfrac{64}{27}$。

故选 D。
\end{answer}

\begin{answer}
\gksolhead{17}{单项选择题}
\gkanswer{答案 B}
设半圆弧半径为 $r$、金属杆 $OM$ 的电阻为 $R$。

过程Ⅰ：$OM$ 从 $OQ$ 逆时针转到 $OS$，扫过的面积为半圆的四分之一（圆心角 $90^{\circ}$），
\[
  \Delta S_{1}=\frac{1}{4}\pi r^{2},\qquad
  q_{1}=\frac{\Delta\Phi_{1}}{R}=\frac{B\Delta S_{1}}{R}=\frac{B\pi r^{2}}{4R}.
\]

过程Ⅱ：$OM$ 固定，磁感应强度由 $B$ 增到 $B'$，回路所围面积为半圆面积
\[
  S_{2}=\frac{1}{2}\pi r^{2},\qquad
  q_{2}=\frac{\Delta\Phi_{2}}{R}=\frac{(B'-B)S_{2}}{R}=\frac{(B'-B)\pi r^{2}}{2R}.
\]

由 $q_{1}=q_{2}$ 得
\[
  \frac{B\pi r^{2}}{4R}=\frac{(B'-B)\pi r^{2}}{2R}
  \;\Longrightarrow\;
  \frac{B}{4}=\frac{B'-B}{2}
  \;\Longrightarrow\;
  B'=\frac{3}{2}B,
\]
故 $\dfrac{B'}{B}=\dfrac{3}{2}$。

故选 B。
\end{answer}

\begin{answer}
\gksolhead{18}{单项选择题}
\gkanswer{答案 C}
小球受到的恒力 $F=mg$ 水平向右。从 $a$ 到 $b$（水平段长 $2R$），该力做功
\[
  W_{1}=F\cdot 2R=2mgR .
\]
到 $b$ 点时的动能为 $\frac{1}{2}mv_{b}^{2}=2mgR$，即 $v_{b}^{2}=4gR$。

在圆弧段 $bc$ 上，恒力 $F$ 与重力 $mg$ 的合力大小为 $\sqrt{2}mg$、方向与水平成 $45^{\circ}$ 向下。小球沿圆弧上升至 $c$ 点时，恒力做功 $W_{2}=F\cdot R=mgR$，重力做功 $-mgR$，两者恰好抵消，故 $c$ 点速率仍满足 $v_{c}^{2}=v_{b}^{2}=4gR$，即 $v_{c}=2\sqrt{gR}$。

离开 $c$ 点后小球做斜抛运动。以 $c$ 为原点，竖直向上为 $y$ 轴、水平向右为 $x$ 轴，初速度 $v_{c}=2\sqrt{gR}$ 竖直向上。竖直方向 $y=v_{c}t-\frac{1}{2}gt^{2}$，到达最高点时 $v_{y}=v_{c}-gt=0$，得
\[
  t=\frac{v_{c}}{g}=2\sqrt{\frac{R}{g}},\qquad
  h_{\max}=y(t)=\frac{v_{c}^{2}}{2g}=2R .
\]
在斜抛阶段恒力（水平向右，大小 $mg$）继续做功，水平位移
\[
  x=\frac{1}{2}a_{x}t^{2}=\frac{1}{2}g\left(2\sqrt{\frac{R}{g}}\right)^{2}=2R,
\]
故斜抛阶段恒力做功 $W_{3}=F\cdot x=mg\cdot 2R=2mgR$。

从 $a$ 到轨迹最高点，恒力做功总和
\[
  W=W_{1}+W_{2}+W_{3}=2mgR+mgR+2mgR=5mgR .
\]
除重力外只有恒力 $F$ 做功，故机械能增量等于恒力做的功，即 $\Delta E=5mgR$。

故选 C。
\end{answer}

\begin{answer}
\gksolhead{19}{多项选择题}
\gkanswer{答案 AD}
开关闭合的瞬间，与电源相连的线圈中电流由零增大，铁芯中磁通量增大，穿过另一线圈的磁通量随之增大，由楞次定律，另一线圈中产生感应电流以阻碍磁通量增加。该感应电流沿远处直导线流动，直导线在小磁针处产生的磁场使 N 极发生偏转；按楞次定律「阻碍增加」的方向判定，直导线电流在小磁针处的磁场垂直纸面向里，N 极朝纸面向里转动，A 正确。

开关闭合并保持一段时间后，磁通量不再变化，另一线圈中无感应电流，直导线无电流、无磁场，小磁针恢复原指向（沿地磁场南北方向），不会稳定指向纸面内或外，B、C 错误。

开关断开瞬间，磁通量迅速减小，感应电流方向与闭合瞬间相反，小磁针处的磁场垂直纸面向外，N 极朝纸面向外转动，D 正确。

故选 AD。
\end{answer}

\begin{answer}
\gksolhead{20}{多项选择题}
\gkanswer{答案 BC}
两颗中子星绕连线上某点做匀速圆周运动，周期相同。设质量分别为 $m_{1}$、$m_{2}$，轨道半径分别为 $r_{1}$、$r_{2}$，角速度 $\omega=2\pi n$（$n=12\,\si{\per\second}$），间距 $L=r_{1}+r_{2}=400\,\si{\kilo\metre}$。

由万有引力提供向心力：
\[
  \frac{Gm_{1}m_{2}}{L^{2}}=m_{1}\omega^{2}r_{1}=m_{2}\omega^{2}r_{2}.
\]
两式相加并利用 $r_{1}+r_{2}=L$：
\[
  \frac{G(m_{1}+m_{2})}{L^{2}}=\omega^{2}L
  \;\Longrightarrow\;
  m_{1}+m_{2}=\frac{\omega^{2}L^{3}}{G}.
\]
已知 $n$、$L$、$G$ 即可求**质量之和**，故 B 正确；质量之积 $\dfrac{\omega^{2}r_{1}r_{2}L}{G}$ 需先知道 $r_{1}$、$r_{2}$ 的分配，无法单独确定，A 错误。

两星速率之和
\[
  v_{1}+v_{2}=\omega r_{1}+\omega r_{2}=\omega L=2\pi nL,
\]
可由已知量直接算出，故 C 正确。

自转角速度与公转无关，题给条件无法确定，D 错误。

故选 BC。
\end{answer}

\begin{answer}
\gksolhead{21}{多项选择题}
\gkanswer{答案 AB}
相邻等势面间距相等，故相邻两等势面间的电势差相等，设为 $\Delta\varphi$。电子从 $a$ 到 $d$ 克服电场力做功 $6\,\si{eV}$，经过 3 个间距，故
\[
  3e\Delta\varphi=6\,\si{eV}\;\Longrightarrow\;\Delta\varphi=2\,\si{\volt}.
\]
又已知 $\varphi_{b}=2\,\si{\volt}$。电子带负电，从 $a$ 到 $d$ 克服电场力做功说明电场力做负功，即电势沿 $a\to d$ 方向降低，故 $\varphi_{a}>\varphi_{b}>\varphi_{c}>\varphi_{d}$。

由 $\varphi_{b}=2\,\si{\volt}$、相邻差 $2\,\si{\volt}$ 得 $\varphi_{c}=0$，A 正确。

电子到 $d$ 时动能 $E_{\mathrm{k}d}=10-6=4\,\si{eV}$。若继续向 $f$ 运动，还需克服电场力做功 $2\times2=4\,\si{eV}$，恰好耗尽动能，即电子最多到达 $f$ 面（速度减为零）。但电子在 $a$ 时的速度方向未必沿电场反方向，若入射方向偏离该方向，则到达 $f$ 面前就已折返，因此**可能到达不了** $f$ 面，B 正确。

电子在 $d$ 点的电势能 $E_{\mathrm{p}d}=-e\varphi_{d}$，而 $\varphi_{d}=\varphi_{b}-2\Delta\varphi=2-4=-2\,\si{\volt}$，故 $E_{\mathrm{p}d}=+2\,\si{eV}\neq 4\,\si{eV}$，C 错误。

电子在 $b$ 点动能 $E_{\mathrm{k}b}=10-2=8\,\si{eV}$，在 $d$ 点 $E_{\mathrm{k}d}=4\,\si{eV}$，速率之比 $\sqrt{8/4}=\sqrt{2}$，不是 2 倍，D 错误。

故选 AB。
\end{answer}

\begin{answer}
\gksolhead{22}{实验题}
\gkanswer{答案 3.775；53.7}
游标卡尺（主尺最小分度 \SI{1}{\milli\metre}，游标 20 分度，精度 \SI{0.05}{\milli\metre}）：图（b）主尺读数为 \SI{3.7}{\centi\metre}，游标第 15 条刻线与主尺刻线对齐，故
\[
  l=3.7\,\si{\centi\metre}+15\times0.005\,\si{\centi\metre}=3.775\,\si{\centi\metre}.
\]
（读作 \SI{3.775}{\centi\metre} 或按精度记作 $3.775$。）

弹簧伸长量
\[
  \Delta x=l-l_{0}=3.775\,\si{\centi\metre}-1.950\,\si{\centi\metre}=1.825\,\si{\centi\metre}=\SI{1.825e-2}{\metre}.
\]
由 $mg=k\Delta x$ 得
\[
  k=\frac{mg}{\Delta x}=\frac{0.100\times9.80}{1.825\times10^{-2}}\,\si{\newton\per\metre}\approx\SI{53.7}{\newton\per\metre}.
\]
\end{answer}

\begin{answer}
\gksolhead{23}{实验题}
\gkanswer{答案 (1) $b$；(2) 见图线；(3) 450；(4) 620.0，33.0}
(1) 为保护电路，闭合 $\mathrm{S}_{1}$ 前应使滑动变阻器接入电阻最大，故滑片移到 $b$ 端。

(2) 按表格数据描点，$R_{2}$ 随 $t$ 增大而单调减小，连线得到一条平滑下降的曲线（见图（b））。

(3) 由 $R_{2}$--$t$ 图线读取 $t=44.0\,^{\circ}\mathrm{C}$ 时的 $R_{2}$ 值约为 $450\,\si{\ohm}$。因 $\mathrm{S}_{2}$ 分别接 1、2 时电压表读数相同，说明两次回路电流相同，故 $R_{T}=R_{2}$，即 $R_{T}\approx450\,\si{\ohm}$。

(4) 图（c）电阻箱读数为 $6\times100+2\times10+0\times1+0\times0.1=620.0\,\si{\ohm}$。由 $R_{2}$--$t$ 图线反查 $R_{2}=620.0\,\si{\ohm}$ 对应的温度约为 $33.0\,^{\circ}\mathrm{C}$。
\end{answer}

\begin{answer}
\gksolhead{24}{解答题}
\gkanswer{答案 (1) $t=\dfrac{1}{g}\sqrt{\dfrac{2E}{m}}$；(2) $h=\dfrac{2E}{mg}$}
(1) 设烟花弹上升的初速度为 $v_{0}$，由题给条件
\[
  E=\frac{1}{2}mv_{0}^{2}. \tag{①}
\]
设从地面上升到火药爆炸所用时间为 $t$，由运动学公式（上升至速度为零）
\[
  0-v_{0}=-gt. \tag{②}
\]
联立①②得
\[
  t=\frac{1}{g}\sqrt{\frac{2E}{m}}.
\]

(2) 设爆炸时烟花弹距地面高度为 $h_{1}$，由机械能守恒
\[
  E=mgh_{1}. \tag{③}
\]
爆炸后上、下两部分质量均为 $\dfrac{m}{2}$，设炸后瞬间速度分别为 $v_{1}$、$v_{2}$。由题给条件（动能之和也为 $E$）和动量守恒：
\[
  \frac{1}{4}mv_{1}^{2}+\frac{1}{4}mv_{2}^{2}=E, \tag{④}
\]
\[
  \frac{1}{2}mv_{1}+\frac{1}{2}mv_{2}=0. \tag{⑤}
\]
由⑤知两部分速度方向相反，向上运动的部分做竖直上抛。设其继续上升的高度为 $h_{2}$，由机械能守恒
\[
  \frac{1}{4}mv_{1}^{2}=\frac{1}{2}mg h_{2}\cdot\frac{1}{2}\quad\Longleftrightarrow\quad
  \frac{1}{4}mv_{1}^{2}=\frac{1}{4}mg\cdot 2h_{2}. \tag{⑥}
\]
由④⑤解得 $\dfrac{1}{4}mv_{1}^{2}=\dfrac{E}{2}$，代入⑥得 $h_{2}=\dfrac{E}{mg}$。

于是上部分距地面的最大高度
\[
  h=h_{1}+h_{2}=\frac{E}{mg}+\frac{E}{mg}=\frac{2E}{mg}.
\]
\end{answer}

\begin{answer}
\gksolhead{25}{解答题}
\gkanswer{答案 (1) $s_{1}=\dfrac{2\sqrt{3}}{3}h$；(2) $B=\sqrt{\dfrac{6mE}{qh}}$；(3) $s_{2}'=\dfrac{2\sqrt{3}}{3}\left(\sqrt{2}-1\right)h$}
(1) ${}_{1}^{1}\mathrm{H}$ 在电场中做类平抛运动，在磁场中做匀速圆周运动。设电场中加速度为 $a_{1}$、初速度为 $v_{1}$、运动时间为 $t_{1}$，第一次进入磁场的位置到 $O$ 的距离为 $s_{1}$。由运动学公式
\[
  s_{1}=v_{1}t_{1}, \tag{①}
\]
\[
  h=\frac{1}{2}a_{1}t_{1}^{2}. \tag{②}
\]
进入磁场时速度方向与 $x$ 轴正方向夹角 $\theta_{1}=60^{\circ}$，故
\[
  a_{1}t_{1}=v_{1}\tan 60^{\circ}=\sqrt{3}v_{1}. \tag{③}
\]
由牛顿第二定律
\[
  qE=ma_{1}. \tag{④}
\]
联立①②③④解得
\[
  s_{1}=\frac{2\sqrt{3}}{3}h.
\]

(2) 进入磁场时速度大小
\[
  v_{1}'=\sqrt{v_{1}^{2}+(a_{1}t_{1})^{2}}=2v_{1}.
\]
在磁场中做圆周运动，洛伦兹力提供向心力：
\[
  qv_{1}'B=\frac{mv_{1}'^{2}}{R_{1}}. \tag{⑤}
\]
轨迹自进入点（在 $y=0$ 上、距 $O$ 为 $s_{1}$）到 $O$ 点射出，弦长等于 $s_{1}$，由几何关系
\[
  s_{1}=2R_{1}\sin 60^{\circ}=\sqrt{3}R_{1}. \tag{⑥}
\]
联立①②③④⑤⑥，并利用 $E_{\mathrm{k}}=\dfrac{1}{2}m v_{1}^{2}$ 与 $h$ 的关系，解得
\[
  B=\sqrt{\frac{6mE}{qh}}.
\]

(3) ${}_{1}^{2}\mathrm{H}$ 的质量为 $2m$，电荷量为 $q$。以相同动能射出，则 $\dfrac{1}{2}(2m)v_{2}^{2}=E$，故 $v_{2}=\dfrac{v_{1}}{\sqrt{2}}$。

电场中加速度
\[
  qE=2ma_{2}\;\Longrightarrow\;a_{2}=\frac{a_{1}}{2}.
\]
设其在电场中运动时间为 $t_{2}$，由 $h=\dfrac{1}{2}a_{2}t_{2}^{2}$ 得 $t_{2}=\sqrt{2}\,t_{1}$。于是
\[
  s_{2}=v_{2}t_{2}=\frac{v_{1}}{\sqrt{2}}\cdot\sqrt{2}\,t_{1}=s_{1},
\]
\[
  a_{2}t_{2}=\frac{a_{1}}{2}\cdot\sqrt{2}t_{1}=\frac{\sqrt{2}}{2}\,a_{1}t_{1}
  \;\Longrightarrow\;
  \tan\theta_{2}=\frac{a_{2}t_{2}}{v_{2}}=\frac{(\sqrt{2}/2)a_{1}t_{1}}{v_{1}/\sqrt{2}}=a_{1}t_{1}\cdot\frac{1}{v_{1}}\cdot\sqrt{2}\cdot\frac{\sqrt{2}}{2}=\tan\theta_{1},
\]
即 $\theta_{2}=\theta_{1}=60^{\circ}$。进入磁场时速度
\[
  v_{2}'=\sqrt{v_{2}^{2}+(a_{2}t_{2})^{2}}=\frac{\sqrt{2}}{2}v_{1}'.
\]
磁场中半径
\[
  R_{2}=\frac{2mv_{2}'}{qB}=\frac{2m}{qB}\cdot\frac{\sqrt{2}}{2}v_{1}'=\sqrt{2}R_{1}.
\]
射出点与进入点的距离
\[
  s_{2}'=2R_{2}\sin 60^{\circ}=\sqrt{2}\cdot 2R_{1}\sin60^{\circ}=\sqrt{2}s_{1}.
\]
由于两核进入磁场的点都在 $y=0$ 上、距 $O$ 同为 $s_{1}=s_{2}$，而 ${}_{1}^{2}\mathrm{H}$ 的运动方向使它在 $O$ 的另一侧射出，故其第一次离开磁场的位置到 $O$ 的距离为
\[
  s_{2}'-s_{2}=\left(\sqrt{2}-1\right)s_{1}=\frac{2\sqrt{3}}{3}\left(\sqrt{2}-1\right)h.
\]
\end{answer}

\begin{answer}
\gksolhead{33}{选考题（选修 3-3）}
\gkanswer{答案 (1) BDE；(2) $m_{\text{液}}=\dfrac{15p_{0}S}{26g}$}
(1) 由图（a），过程①为等容升压升温（$V$ 不变、$p$ 增大、$T$ 升高），A 错误。过程②体积增大，气体对外界做正功，B 正确。过程④中温度降低，由热力学第一定律 $\Delta U=Q+W$，体积不变则 $W=0$，$\Delta U<0$ 故 $Q<0$，即气体向外界放热，C 错误。状态 $c$、$d$ 温度相同，理想气体内能只由温度决定，故内能相等，D 正确。由 $\dfrac{pV}{T}=\text{常量}$，$T_{b}=T_{d}$ 而 $V_{b}<V_{d}$，故 $p_{b}>p_{d}$，即状态 $d$ 的压强比状态 $b$ 的小，E 正确。

(2) 设活塞再次平衡后，上方气体体积为 $V_{1}$、压强为 $p_{1}$，下方气体体积为 $V_{2}$、压强为 $p_{2}$。活塞下移过程中上、下方气体温度均不变，由玻意耳定律
\[
  p_{0}\cdot\frac{V}{2}=p_{1}V_{1}, \tag{①}
\]
\[
  p_{0}\cdot\frac{V}{2}=p_{2}V_{2}. \tag{②}
\]
由已知条件：流入液体体积为 $\dfrac{V}{8}$，下方气体体积减小了 $\dfrac{V}{6}$，故
\[
  V_{2}=\frac{V}{2}-\frac{V}{6}=\frac{V}{3}.
\]
上方空间被液体占据一部分，剩余气体体积
\[
  V_{1}=\frac{V}{2}-\frac{V}{8}+\frac{V}{6}=\frac{13V}{24}.
\]
由①②得
\[
  p_{1}=\frac{p_{0}V/2}{13V/24}=\frac{12p_{0}}{13},\qquad
  p_{2}=\frac{p_{0}V/2}{V/3}=\frac{3p_{0}}{2}.
\]
对活塞受力分析：活塞上方受气体压力 $p_{1}S$ 与液体重量 $m_{\text{液}}g$，下方受气体压力 $p_{2}S$（活塞质量、体积不计），故
\[
  p_{2}S=p_{1}S+m_{\text{液}}g
  \;\Longrightarrow\;
  m_{\text{液}}=\frac{(p_{2}-p_{1})S}{g}.
\]
代入得
\[
  m_{\text{液}}=\frac{\left(\dfrac{3}{2}-\dfrac{12}{13}\right)p_{0}S}{g}
  =\frac{\dfrac{39-24}{26}\,p_{0}S}{g}
  =\frac{15p_{0}S}{26g}.
\]
\end{answer}

\begin{answer}
\gksolhead{34}{选考题（选修 3-4）}
\gkanswer{答案 (1) $\sqrt{3}$，大于；(2)(i) $v=\SI{18}{\centi\metre\per\second}$，沿 $x$ 轴负方向；(ii) $x_{Q}=\SI{9}{\centi\metre}$}
(1) 红光垂直 $AB$ 面射入，在 $AB$ 面上不发生偏折；射到 $AC$ 面时，入射角等于棱镜顶角 $\angle A=30^{\circ}$，折射角 $r=60^{\circ}$，故
\[
  n=\frac{\sin r}{\sin i}=\frac{\sin60^{\circ}}{\sin30^{\circ}}=\sqrt{3}.
\]
蓝光折射率大于红光，同一入射角下折射角更大，故蓝光射出时的折射角**大于** $60^{\circ}$。

(2)(i) 由图（a）读出波长
\[
  \lambda=\SI{36}{\centi\metre}.
\]
由图（b）读出周期
\[
  T=\SI{2}{\second}.
\]
波速
\[
  v=\frac{\lambda}{T}=\frac{36}{2}\,\si{\centi\metre\per\second}=\SI{18}{\centi\metre\per\second}.
\]
由图（b）知，当 $t=\dfrac{1}{3}\,\si{\second}$ 时 $Q$ 点向上运动；结合图（a）可知波沿**负方向**传播。

(ii) 设质点 $P$、$Q$ 平衡位置的 $x$ 坐标分别为 $x_{P}$、$x_{Q}$。由图（a）知，$x=0$ 处
\[
  y=-\frac{A}{2}=A\sin(-30^{\circ}). \tag{④}
\]
因此
\[
  x_{P}=\frac{30^{\circ}}{360^{\circ}}\lambda=\SI{3}{\centi\metre}.
\]
由图（b）知，在 $t=0$ 时 $Q$ 点处于平衡位置；经
\[
  \Delta t=\frac{1}{3}\,\si{\second}
\]
其振动状态沿 $x$ 轴负方向传播至 $P$ 点处，由此及波速得
\[
  x_{Q}-x_{P}=v\Delta t=18\times\frac{1}{3}\,\si{\centi\metre}=\SI{6}{\centi\metre}. \tag{⑤}
\]
由④⑤得质点 $Q$ 的平衡位置的 $x$ 坐标
\[
  x_{Q}=\SI{9}{\centi\metre}.
\]
\end{answer}
